\section{Conclusion}
\label{sec:conclusion}

We set out to measure, rather than derive, the tension between group and
individual fairness. To do so we built a differentiable composite loss whose
parameter $\beta$ interpolates continuously between a soft demographic-parity
term and a soft generalized-entropy term, and swept it across a $7\times7$
$(\alpha,\beta)$ grid, three datasets, two model families and ten seeds:
$2{,}940$ trained models with paired significance testing throughout.

Three findings stand out. First, the two objectives exhibit a statistically
significant \emph{double dissociation}: the group term reduces demographic
disparity by $78\%$--$89\%$ while significantly worsening benefit inequality,
and the individual term significantly improves benefit inequality while
leaving disparity unchanged or worse. The effects are strongly asymmetric in
magnitude, because equalizing benefit is bounded by the achievable prediction
quality in a way that equalizing selection rates is not. Second, on Adult the
conflict extends within group fairness: driving demographic parity to $0.009$
raises equalized odds by a factor of $3.3$, an empirical instance of a known
impossibility, with the exchange rate quantified. Third, where a baseline
already exhibits little disparity, the intervention has nothing to act on, and
where the dataset is small, validation-based selection cannot reliably find
the configurations that would act---on German Credit the tuned model
reproduced the baseline disparity while a pre-specified configuration reduced
it by $48\%$ at no accuracy cost.

The composite loss reaches group-fairness levels comparable to established
pre-, in- and post-processing methods without surpassing them. That was not
the objective; the objective was a measurement instrument credible on the axis
it shares with existing tools, and one that additionally traverses the axis
they do not.

The practical consequence is that $\beta$ should not be tuned. There is no
metric that ranks its endpoints, because they optimize different and partly
opposed definitions of the same word. It is a policy parameter, and the
contribution of a trade-off map is to attach a measured price to each setting
so that the choice is made deliberately rather than inherited from a library
default.

Future work should extend the map to intersectional subgroups, replace the
all-groups equalized odds statistic with a largest-two-groups variant on
COMPAS, test the composite loss under mini-batch training at larger scale, and
characterise more precisely the sample-size floor below which fairness
model selection ceases to function.
